% GENERATED FILE. Edit canonical lessons, then run npm run generate:books.
% canonical-source-hash: fnv1a64:5b1589fdf5cf2a01
% canonical-lessons: ZH-C401-qiantu, ZH-C401-gaihang, ZH-C401-chengzhang, ZH-C401-bianlao, ZH-C401-yiqie, ZH-R401-review-397-401, ZH-R401-review-final-401

\chapter{Plans for the Future}
\label{ch:zh-plans-for-the-future}
\hlchaptermodality{401}

\begin{chapteropening}
\textbf{By the end of this chapter:} \emph{I can talk about prospects, changing career, growing up, getting older and everything being fine.}

Complete the last lesson of chapter 401: I can talk about prospects, changing career, growing up, getting older and everything being fine.
\end{chapteropening}

\section[qiántú]{qiántú — prospects; future (of a person or job)}

\label{lesson:ZH-C401-qiantu}



\begin{quote}
Before the new one: say the Mandarin for originally planned, then the Mandarin for settled and secure.
\end{quote}

\subsection*{You'll want to know: qiántú}
\emph{qiántú} — \textquotedblleft{}prospects; future (of a person or job)\textquotedblright{}.

\begin{quote}
\emph{wǒ juéde zhège gōngzuò hěn yǒu qiántú} — I think this job has good prospects.
\end{quote}

Its characters come later; for now it is a word for the ear.

\begin{cousinweb}
\emph{qiántú} is built from \zh{前} 'ahead' and \emph{tú} 'a road', both inherited from Middle Chinese, \emph{dzen} and \emph{du} (Guangyun, in Baxter's transcription), which glosses the second 'a way': the road ahead.
\end{cousinweb}

\begin{grammarlens}[title={What you've built}]
One more word for saying what would happen.
\end{grammarlens}

\subsection*{Your turn}
\begin{itemize}
\raggedright
  \item \emph{Say it:} \emph{qiántú}
  \item \emph{Say it:} \emph{qiántú}, once more
  \item \emph{Say it:} \emph{wǒ juéde zhège gōngzuò hěn yǒu qiántú} — I think this job has good prospects.
  \item \emph{Recall:} say the Mandarin for originally planned, then the Mandarin for settled and secure, then say \emph{qiántú} again
\end{itemize}

\begin{tcolorbox}[breakable,colback=teal!4,colframe=teal!35!black,title={Before you move on}]
What does qiántú mean? (\textquotedblleft{}Prospects; future (of a person or job)\textquotedblright{}.) Say it once more.
\end{tcolorbox}
\section[gǎiháng]{gǎiháng — to change careers}

\label{lesson:ZH-C401-gaihang}



\begin{quote}
Before the new one: say the Mandarin for settled and secure, then the Mandarin for future prospects.
\end{quote}

\subsection*{You'll want to know: gǎiháng}
\emph{gǎiháng} — \textquotedblleft{}to change careers\textquotedblright{}.

\begin{quote}
\emph{yàoshi bù xǐhuan zhège gōngzuò wǒ jiù gǎiháng} — If I don't like this job, I'll change careers.
\end{quote}

Its characters come later; for now it is a word for the ear.

\begin{cousinweb}
\emph{gǎiháng} is built from \zh{改} 'to alter' and \emph{háng} 'a line', both inherited from Middle Chinese, \emph{kojX} and \emph{hang} (Guangyun, in Baxter's transcription), which glosses the second 'a rank' and 'a row'. A line of work is a later sense.
\end{cousinweb}

\begin{grammarlens}[title={What you've built}]
One more word for saying what would happen.
\end{grammarlens}

\subsection*{Your turn}
\begin{itemize}
\raggedright
  \item \emph{Say it:} \emph{gǎiháng}
  \item \emph{Say it:} \emph{gǎiháng}, once more
  \item \emph{Say it:} \emph{yàoshi bù xǐhuan zhège gōngzuò wǒ jiù gǎiháng} — If I don't like this job, I'll change careers.
  \item \emph{Recall:} say the Mandarin for settled and secure, then the Mandarin for future prospects, then say \emph{gǎiháng} again
\end{itemize}

\begin{tcolorbox}[breakable,colback=teal!4,colframe=teal!35!black,title={Before you move on}]
What does gǎiháng mean? (\textquotedblleft{}To change careers\textquotedblright{}.) Say it once more.
\end{tcolorbox}
\section[\texorpdfstring{\zh{成长} (chéngzhǎng)}{成长 (chéngzhǎng)}]{\zh{成长} (chéngzhǎng) — to mature (as a person)}

\label{lesson:ZH-C401-chengzhang}



\begin{quote}
Before the new one: say the Mandarin for future prospects, then the Mandarin for to change careers.
\end{quote}

\subsection*{You'll want to know: \zh{成长}}
\textbf{\zh{成长}} — \emph{chéngzhǎng} — \textquotedblleft{}to mature (as a person)\textquotedblright{}.

\begin{quote}
\textbf{\zh{孩子成长得很快}} \emph{háizi chéngzhǎng de hěn kuài} — Children grow up fast.
\end{quote}

\begin{cousinweb}
\textbf{\zh{成长}} (\emph{chéngzhǎng}) is built from \zh{成} 'to complete' and \zh{长} 'to grow', both inherited from Middle Chinese, \emph{dzyeng} and \emph{trjangX} (Guangyun, in Baxter's transcription), which glosses the second 'big'.
\end{cousinweb}

\begin{grammarlens}[title={What you've built}]
One more word for saying what would happen.
\end{grammarlens}

\subsection*{Your turn}
\begin{itemize}
\raggedright
  \item \emph{Say it:} \emph{chéngzhǎng}
  \item \emph{Say it:} \emph{chéngzhǎng}, once more
  \item \emph{Say it:} \emph{háizi chéngzhǎng de hěn kuài} — Children grow up fast.
  \item \emph{Recall:} say the Mandarin for future prospects, then the Mandarin for to change careers, then say \emph{chéngzhǎng} again
\end{itemize}

\begin{tcolorbox}[breakable,colback=teal!4,colframe=teal!35!black,title={Before you move on}]
What does \zh{成长} mean? (\textquotedblleft{}To mature (as a person)\textquotedblright{}.) Say it once more.
\end{tcolorbox}
\section[\texorpdfstring{\zh{变老} (biànlǎo)}{变老 (biànlǎo)}]{\zh{变老} (biànlǎo) — to get old}

\label{lesson:ZH-C401-bianlao}



\begin{quote}
Before the new one: say the Mandarin for to change careers, then the Mandarin for to mature (as a person).
\end{quote}

\subsection*{You'll want to know: \zh{变老}}
\textbf{\zh{变老}} — \emph{biànlǎo} — \textquotedblleft{}to get old\textquotedblright{}.

\begin{quote}
\textbf{\zh{人都会变老}} \emph{rén dōu huì biànlǎo} — Everyone gets old.
\end{quote}

\begin{cousinweb}
\textbf{\zh{变老}} (\emph{biànlǎo}) is built from \zh{变} 'to change' and \zh{老} 'old', both inherited from Middle Chinese, \emph{pjenH} and \emph{lawX} (Guangyun, in Baxter's transcription), which glosses the second 'aged'.
\end{cousinweb}

\begin{grammarlens}[title={What you've built}]
One more word for saying what would happen.
\end{grammarlens}

\subsection*{Your turn}
\begin{itemize}
\raggedright
  \item \emph{Say it:} \emph{biànlǎo}
  \item \emph{Say it:} \emph{biànlǎo}, once more
  \item \emph{Say it:} \emph{rén dōu huì biànlǎo} — Everyone gets old.
  \item \emph{Recall:} say the Mandarin for to change careers, then the Mandarin for to mature (as a person), then say \emph{biànlǎo} again
\end{itemize}

\begin{tcolorbox}[breakable,colback=teal!4,colframe=teal!35!black,title={Before you move on}]
What does \zh{变老} mean? (\textquotedblleft{}To get old\textquotedblright{}.) Say it once more.
\end{tcolorbox}
\section[yíqiè]{yíqiè — everything; all}

\label{lesson:ZH-C401-yiqie}



\begin{quote}
Before the new one: say the Mandarin for to mature (as a person), then the Mandarin for to get older.
\end{quote}

\subsection*{You'll want to know: yíqiè}
\emph{yíqiè} — \textquotedblleft{}everything; all\textquotedblright{}.

\begin{quote}
\emph{zhǐyào nǐ zài yíqiè dōu hǎo} — As long as you're here, everything's fine.
\end{quote}

Its characters come later; for now it is a word for the ear.

Tone: \zh{一} is said yí here, rising, because qiè after it is a falling tone.

\begin{cousinweb}
\emph{yíqiè} is built from \zh{一} 'one' and \emph{qiè}, both inherited from Middle Chinese, \emph{'jit} and \emph{tshejH} (Guangyun, in Baxter's transcription), which glosses the second, in this tone, 'the many': all of them as one.
\end{cousinweb}

\begin{grammarlens}[title={What you've built}]
One more word for saying what would happen.
\end{grammarlens}

\subsection*{Your turn}
\begin{itemize}
\raggedright
  \item \emph{Say it:} \emph{yíqiè}
  \item \emph{Say it:} \emph{yíqiè}, once more
  \item \emph{Say it:} \emph{zhǐyào nǐ zài yíqiè dōu hǎo} — As long as you're here, everything's fine.
  \item \emph{Recall:} say the Mandarin for to mature (as a person), then the Mandarin for to get older, then say \emph{yíqiè} again
\end{itemize}

\begin{tcolorbox}[breakable,colback=teal!4,colframe=teal!35!black,title={Before you move on}]
What does yíqiè mean? (\textquotedblleft{}Everything; all\textquotedblright{}.) Say it once more.
\end{tcolorbox}
\section[(dialogue)]{Looking back: from Surely Not to Plans for the Future}

\label{lesson:ZH-R401-review-397-401}



\begin{quote}
Say the Mandarin for to stand someone up and surely not (in a question).
\end{quote}

\subsection*{Your turn}
\begin{itemize}
\raggedright
  \item \emph{Say it:} to stand someone up, surely not (in a question), hard to come by, in an awkward spot, an excuse (a poor reason)
  \item \emph{Say it:} your attitude, the difference, a region of a country, ready-made, to hinge on
  \item \emph{Say it:} to demand (of someone), only if, as long as, to meet (a requirement), to marry and settle down
  \item \emph{Say it:} to adjust, to extend (in time), to shorten, originally planned, settled and secure
  \item \emph{Say it:} future prospects, to change careers, to mature (as a person), to get older, everything
  \item \emph{Read it:} \zh{难} \zh{只} — each one aloud, then a word that holds it
\end{itemize}

\begin{tcolorbox}[breakable,colback=teal!4,colframe=teal!35!black,title={Before you move on}]
To stand someone up? (\emph{fàng gēzi}.) Surely not (in a question)? (\textbf{\zh{难道}}, \emph{nándào}.) Hard to come by? (\textbf{\zh{难得}}, \emph{nándé}.) In an awkward spot? (\textbf{\zh{为难}}, \emph{wéinán}.) An excuse (a poor reason)? (\emph{jièkǒu}.) Your attitude? (\emph{tàidu}.) The difference? (\emph{qūbié}.) A region of a country? (\textbf{\zh{地区}}, \emph{dìqū}.) Ready-made? (\textbf{\zh{现成}}, \emph{xiànchéng}.) To hinge on? (\emph{qǔjuéyú}.) To demand (of someone)? (\emph{yāoqiú}.) Only if? (\textbf{\zh{只有}}, \emph{zhǐyǒu}.) As long as? (\textbf{\zh{只要}}, \emph{zhǐyào}.) To meet (a requirement)? (\emph{mǎnzú}.) To marry and settle down? (\textbf{\zh{成家}}, \emph{chéngjiā}.) To adjust? (\emph{tiáozhěng}.) To extend (in time)? (\emph{yáncháng}.) To shorten? (\emph{suōduǎn}.) Originally planned? (\emph{yuándìng}.) Settled and secure? (\textbf{\zh{安定}}, \emph{āndìng}.) Future prospects? (\emph{qiántú}.) To change careers? (\emph{gǎiháng}.) To mature (as a person)? (\textbf{\zh{成长}}, \emph{chéngzhǎng}.) To get older? (\textbf{\zh{变老}}, \emph{biànlǎo}.) Everything? (\emph{yíqiè}.)
\end{tcolorbox}
\section[(dialogue)]{Plans for the Future — once more}

\label{lesson:ZH-R401-review-final-401}



\begin{quote}
Say the Mandarin for future prospects and to change careers.
\end{quote}

\subsection*{Your turn}
\begin{itemize}
\raggedright
  \item \emph{Say it:} future prospects, to change careers, to mature (as a person), to get older, everything
\end{itemize}

\begin{tcolorbox}[breakable,colback=teal!4,colframe=teal!35!black,title={Before you move on}]
Future prospects? (\emph{qiántú}.) To change careers? (\emph{gǎiháng}.) To mature (as a person)? (\textbf{\zh{成长}}, \emph{chéngzhǎng}.) To get older? (\textbf{\zh{变老}}, \emph{biànlǎo}.) Everything? (\emph{yíqiè}.)
\end{tcolorbox}
